\section{Problem Data Blocks}\label{sec:problem-data-blocks}

A problem file contains one or more problems or cases, which in turn contain one or
more trajectory definitions. Each trajectory definition begins with a
\problemblock{*trajectory} block (\cref{sec:trajectory-data-blocks}) and ends at
the next nontrajectory, nonsegment keyword. The trajectories contain one or more
segments and therefore form the hierarchy introduced in
\cref{fig:problem-file-organization}.

It is important to group data blocks according to their purpose. All blocks defining
a segment must be grouped together, and all blocks defining a trajectory must be
grouped together. The remaining problem-level blocks do not have to be contiguous,
but the problem file is easier to understand if they are grouped rather than scattered
throughout the file. Grouping these blocks also tends to reduce input errors.

\begin{figure}[htbp]
  \centering
  \begin{tikzpicture}[font=\small,node distance=1.4mm]
  \node[draw,minimum width=7.4cm,minimum height=0.55cm] (id) {\texttt{(problem-id)}};
  \node[draw,below=of id,minimum width=7.4cm,minimum height=0.55cm] (p1) {Problem data blocks};
  \node[draw,below=of p1,minimum width=7.4cm,minimum height=0.55cm] (t1) {Trajectory 1 data blocks};
  \node[draw,below=of t1,minimum width=6.5cm,minimum height=0.48cm] (s11) {Segment data blocks};
  \node[draw,below=of s11,minimum width=6.5cm,minimum height=0.48cm] (s12) {Segment data blocks};
  \node[draw,below=of s12,minimum width=7.4cm,minimum height=0.55cm] (t2) {Trajectory 2 data blocks};
  \node[draw,below=of t2,minimum width=6.5cm,minimum height=0.48cm] (s21) {Segment data blocks};
  \node[draw,below=of s21,minimum width=6.5cm,minimum height=0.48cm] (s22) {Segment data blocks};
  \node[draw,below=of s22,minimum width=7.4cm,minimum height=0.55cm] (p2) {Problem data blocks};
  \node[draw,below=of p2,minimum width=7.4cm,minimum height=0.55cm] (end) {\texttt{*end}};
  \draw[decorate,decoration={brace,amplitude=5pt},thick]
    ($(id.north east)+(0.25,0)$) -- node[right=7pt]{Problem file}
    ($(end.south east)+(0.25,0)$);
\end{tikzpicture}

  \caption{Problem and data-block organization.}
  \label{fig:problem-data-block-organization}
\end{figure}

The resulting structure is illustrated in
\cref{fig:problem-data-block-organization}. A problem begins with an identifier in
parentheses. The identifier is often the same as the problem filename, but this is not
required. A group of data blocks containing information that applies to the entire
problem follows.

These initial blocks are not associated with a particular trajectory or segment; they
apply to all trajectories. Examples are the \problemblock{*atmos},
\problemblock{*earth}, and \problemblock{*wind} blocks.
\Cref{tab:problem-data-blocks} gives the complete list. Blocks marked with
\limitedstate{} are best placed at the beginning of the problem. The remaining
problem-level blocks should normally be placed after the trajectory definitions.

After the initial group of problem-level blocks, the blocks defining the first
trajectory are given. Within that trajectory are the blocks defining its segments. A
second trajectory definition follows the first, and this pattern continues for every
trajectory in the problem. The final group of problem-level blocks follows the
trajectories, and the problem ends with \problemblock{*end}.

\begin{table}[htbp]
  \centering
  \caption{Problem data blocks.}
  \label{tab:problem-data-blocks}
  \begin{tabularx}{0.82\linewidth}{@{}>{\ttfamily}lX@{}}
    \toprule
    \normalfont Data block & Description \\
    \midrule
    *atmos\limitedstate     & Atmosphere model \\
    *define\limitedstate    & Problem-level user-defined variables \\
    *earth\limitedstate     & Earth model \\
    *egs                    & EGS database files \\
    *file                   & Problem-level output files \\
    *optimize               & Trajectory optimization \\
    *print                  & Problem-level printout \\
    *radar\limitedstate     & Radar observations \\
    *search                 & Trajectory searches \\
    *summarize              & Summary variables \\
    *survey                 & Trajectory surveys \\
    *title\limitedstate     & Problem title \\
    *units/fmt              & Units and output format \\
    *wind\limitedstate      & Wind model \\
    \bottomrule
  \end{tabularx}

  \smallskip
  \footnotesize\limitedstate\ Best placed before trajectory definitions; the
  remaining blocks normally follow the trajectories.
\end{table}

The blocks at the beginning of a problem do not reference trajectory or segment
information such as trajectory numbers, segment numbers, or user-defined variables.
They contain general information independent of the individual trajectories. For
example, \problemblock{*atmos} and \problemblock{*earth} select the atmosphere
and earth models used by every trajectory.

The blocks at the end usually reference trajectory or segment information that must
already have been defined. For example, a \problemblock{*search} block defines an
objective function that refers to trajectory and segment numbers. An
\problemblock{*egs} block may refer to user-defined variables, which must be
defined before they are used. Placing the \problemblock{*egs} block at the end of
the problem avoids this potential forward-reference error.
